\textbf{COLUNI 2026 -- Questão 36}\par\medskip O samba \textit{O Bonde de São Januário}, de Wilson Batista e Ataulfo Alves, lançado em 1940, teve sua primeira versão vetada pela censura do Estado Novo. A letra original dizia: “O bonde de São Januário leva mais um sócio otário, só eu não vou trabalhar.” Após a intervenção da censura, a canção foi alterada e divulgada com os seguintes versos: “\textit{Quem trabalha é que tem razão, eu digo e não tenho medo de errar… O bonde de São Januário leva mais um operário, sou eu que vou trabalhar.}”

Essa modificação da letra da canção ilustra um dos mecanismos de funcionamento do regime varguista. Sobre isso, é CORRETO afirmar:\begin{enumerate}
\item A censura foi aplicada apenas a músicas estrangeiras, não havendo registros de intervenções em composições brasileiras de caráter popular, o que demonstra tolerância do Estado Novo com a cultura nacional.
\item O samba de um modo geral foi adaptado pela censura do Estado Novo, com a finalidade de difundir uma imagem positiva do trabalho, vinculando-o à disciplina e à ordem social.
\item A alteração da letra do samba \textit{O Bonde de São Januário} refletiu uma escolha autônoma dos compositores, interessados em ampliar o alcance da canção junto ao público trabalhador.
\item O caso de censura da canção de Wilson Batista e Ataulfo Alves revela que, durante o Estado Novo, o governo proibiu sistematicamente qualquer exaltação ao trabalho, considerando essa prática um risco à estabilidade do regime.\end{enumerate}